\documentclass[border=5mm]{standalone}
\usepackage[utf8]{inputenc}
\usepackage[T1]{fontenc}
\usepackage{amsmath}
\usepackage{tikz}
\usepackage{pgfplots}
\pgfplotsset{compat=1.18}
\usetikzlibrary{arrows.meta}
\definecolor{utnblue}{RGB}{0, 51, 102}

% Ej7 resultado: y(t) por tramos en 5 zonas
%   Zona I:   t < 0          -> y = 0
%   Zona II:  0 <= t <= 3     -> y = t - 1 + e^{-t}
%   Zona III: 3 < t <= 5      -> y = (2 e^3 + 1) e^{-t}
%   Zona IV:  5 < t <= 8      -> y = 2 e^{3-t} - (t-6) e^{-5}
%   Zona V:   t > 8           -> y = 0
% Pico en t=3 con y(3) = 2 + e^{-3} ≈ 2.05

\begin{document}
\begin{tikzpicture}[>=Latex, font=\small]

\begin{axis}[
    width=14cm, height=6cm,
    axis lines=center,
    xlabel={$t$},
    ylabel={$y(t)$},
    xlabel style={anchor=north west, at={(axis description cs:1,0)}},
    ylabel style={font=\small, rotate=0, anchor=south east, at={(axis description cs:0,1)}},
    xmin=-1, xmax=9.5,
    ymin=-0.2, ymax=2.5,
    clip=false,
    axis line style={->, thick},
    xtick={0,3,5,8},
    ytick={0, 1, 2.05},
    yticklabels={$0$, $1$, $\;2+e^{-3}$},
    every tick label/.style={font=\footnotesize},
]
    % Zona I
    \addplot[utnblue, very thick, domain=-0.9:0, samples=2] {0};
    % Zona II
    \addplot[utnblue, very thick, domain=0:3, samples=120] {x - 1 + exp(-x)};
    % Zona III
    \addplot[utnblue, very thick, domain=3:5, samples=120]
        {(2*exp(3) + 1) * exp(-x)};
    % Zona IV
    \addplot[utnblue, very thick, domain=5:8, samples=160]
        {2*exp(3-x) - (x-6)*exp(-5)};
    % Zona V
    \addplot[utnblue, very thick, domain=8:9.3, samples=2] {0};

    % puntos de empalme
    \fill[utnblue] (axis cs:0, 0) circle (2pt);
    \fill[utnblue] (axis cs:3, {3 - 1 + exp(-3)}) circle (2pt);
    \fill[utnblue] (axis cs:5, {(2*exp(3) + 1) * exp(-5)}) circle (2pt);
    \fill[utnblue] (axis cs:8, 0) circle (2pt);

    % guías de zona en x = 0,3,5,8
    \draw[dashed, gray!55, thin] (axis cs:3, 0) -- (axis cs:3, 2.15);
    \draw[dashed, gray!55, thin] (axis cs:5, 0) -- (axis cs:5, 0.35);
    \draw[dashed, gray!55, thin] (axis cs:8, 0) -- (axis cs:8, 0.35);
    \draw[dashed, gray!55, thin] (axis cs:0, 2.05) -- (axis cs:3, 2.05);

    % etiquetas de zona
    \node[font=\footnotesize, gray] at (axis cs:-0.5, -0.13) {I};
    \node[font=\footnotesize, gray] at (axis cs:1.5, -0.13)  {II};
    \node[font=\footnotesize, gray] at (axis cs:4, -0.13)    {III};
    \node[font=\footnotesize, gray] at (axis cs:6.5, -0.13)  {IV};
    \node[font=\footnotesize, gray] at (axis cs:8.7, -0.13)  {V};
\end{axis}

\end{tikzpicture}
\end{document}
